% TikZ Diagram: Problem Hierarchy showing Master Problems → Sub-Problems → Methods
\begin{figure}[htbp]
\centering
\begin{tikzpicture}[
    scale=0.85,
    every node/.style={font=\small},
    masterbox/.style={rectangle, draw=primaryblue, very thick, fill=primaryblue!10,
                     rounded corners, minimum width=8cm, minimum height=1.2cm, align=center},
    subbox/.style={rectangle, draw=secondaryblue, thick, fill=secondaryblue!5,
                  rounded corners, minimum width=3cm, minimum height=0.8cm, align=center},
    methodbox/.style={rectangle, draw=accentorange, thick, fill=accentorange!10,
                     rounded corners, minimum width=2.5cm, minimum height=0.7cm, align=center},
    arrow/.style={->, >=stealth, thick}
]

% Title
\node[font=\Large\bfseries] at (0, 11) {Dissertation Problem Hierarchy};

% Master Problem A
\node[masterbox] (MPA) at (-6, 9) {
    \textbf{Master Problem A}\\
    \small Adversarial Resilience in\\
    \small Heterogeneous Networks
};

% Master Problem B
\node[masterbox] (MPB) at (6, 9) {
    \textbf{Master Problem B}\\
    \small Robustness-Accuracy-Privacy\\
    \small Trade-off Optimization
};

% Sub-Problems for Master Problem A
\node[subbox] (SPA1) at (-10, 6.5) {\footnotesize SP-A1\\Continuous-Time\\Temporal Graphs};
\node[subbox] (SPA2) at (-6, 6.5) {\footnotesize SP-A2\\Multi-Granularity\\Embeddings};
\node[subbox] (SPA3) at (-10, 4.5) {\footnotesize SP-A3\\Byzantine-Robust\\Fed. Learning};
\node[subbox] (SPA4) at (-6, 4.5) {\footnotesize SP-A4\\Post-Quantum\\Detection};
\node[subbox] (SPA5) at (-2, 6.5) {\footnotesize SP-A5\\Zero-Shot\\Generalization};
\node[subbox] (SPA6) at (-2, 4.5) {\footnotesize SP-A6\\State Space\\Under Poisoning};
\node[subbox] (SPA7) at (-6, 2.5) {\footnotesize SP-A7\\Game-Theoretic\\Evasion};

% Sub-Problems for Master Problem B
\node[subbox] (SPB1) at (2, 6.5) {\footnotesize SP-B1\\Stochastic\\Adversarial Training};
\node[subbox] (SPB2) at (6, 6.5) {\footnotesize SP-B2\\Differentially Private\\Optimal Transport};
\node[subbox] (SPB3) at (10, 6.5) {\footnotesize SP-B3\\Variational\\Bayesian Attention};
\node[subbox] (SPB4) at (2, 4.5) {\footnotesize SP-B4\\Progressive\\Distillation};
\node[subbox] (SPB5) at (6, 4.5) {\footnotesize SP-B5\\Multi-Objective\\Optimization};
\node[subbox] (SPB6) at (10, 4.5) {\footnotesize SP-B6\\Efficient\\Inference};
\node[subbox] (SPB7) at (6, 2.5) {\footnotesize SP-B7\\Online Learning\\Under Drift};

% Novel Methods
\node[methodbox] (M1) at (-10, 0) {\textbf{CT-TGNN}\\98.3\% accuracy};
\node[methodbox] (M2) at (-6, 0) {\textbf{TripleE-TGNN}\\96.8\% accuracy};
\node[methodbox] (M3) at (-2, 0) {\textbf{FedLLM-API}\\87.1\% w/ Byzantine};
\node[methodbox] (M4) at (2, 0) {\textbf{PQ-IDPS}\\96.2\% PQ detection};
\node[methodbox] (M5) at (6, 0) {\textbf{MambaShield}\\91.4\% under poisoning};
\node[methodbox] (M6) at (10, 0) {\textbf{Stochastic Transformer}\\94.2\% cross-dataset};
\node[methodbox] (M7) at (2, -2) {\textbf{Game-Theoretic}\\94.7\% equilibrium};

% Arrows from Master Problems to Sub-Problems
\draw[arrow, primaryblue] (MPA) -- (SPA1);
\draw[arrow, primaryblue] (MPA) -- (SPA2);
\draw[arrow, primaryblue] (MPA) -- (SPA3);
\draw[arrow, primaryblue] (MPA) -- (SPA4);
\draw[arrow, primaryblue] (MPA) -- (SPA5);
\draw[arrow, primaryblue] (MPA) -- (SPA6);
\draw[arrow, primaryblue] (MPA) -- (SPA7);

\draw[arrow, primaryblue] (MPB) -- (SPB1);
\draw[arrow, primaryblue] (MPB) -- (SPB2);
\draw[arrow, primaryblue] (MPB) -- (SPB3);
\draw[arrow, primaryblue] (MPB) -- (SPB4);
\draw[arrow, primaryblue] (MPB) -- (SPB5);
\draw[arrow, primaryblue] (MPB) -- (SPB6);
\draw[arrow, primaryblue] (MPB) -- (SPB7);

% Arrows from Sub-Problems to Methods (Primary mappings)
\draw[arrow, accentorange] (SPA1) -- (M1);
\draw[arrow, accentorange] (SPA2) -- (M2);
\draw[arrow, accentorange] (SPA3) -- (M3);
\draw[arrow, accentorange] (SPA4) -- (M4);
\draw[arrow, accentorange] (SPA5) -- (M3);
\draw[arrow, accentorange] (SPA6) -- (M5);
\draw[arrow, accentorange] (SPA7) -- (M7);

\draw[arrow, accentorange] (SPB1) -- (M4);
\draw[arrow, accentorange] (SPB1) -- (M6);
\draw[arrow, accentorange] (SPB2) -- (M3);
\draw[arrow, accentorange] (SPB3) -- (M6);
\draw[arrow, accentorange] (SPB4) -- (M5);
\draw[arrow, accentorange] (SPB6) -- (M1);
\draw[arrow, accentorange] (SPB6) -- (M2);
\draw[arrow, accentorange] (SPB6) -- (M5);
\draw[arrow, accentorange] (SPB7) -- (M7);

% Legend
\node[font=\footnotesize, align=left] at (-6, -3.5) {
    \textbf{Legend:}\\
    \textcolor{primaryblue}{\rule{0.5cm}{0.1cm}} Master Problems\\
    \textcolor{secondaryblue}{\rule{0.5cm}{0.1cm}} Sub-Problems\\
    \textcolor{accentorange}{\rule{0.5cm}{0.1cm}} Novel Methods
};

% Mathematical Frameworks annotation
\node[font=\footnotesize, align=left, draw=darkgreen, thick, rounded corners, fill=darkgreen!5] at (6, -3.5) {
    \textbf{Mathematical Frameworks:}\\
    $\bullet$ Neural ODEs (M1)\\
    $\bullet$ Graph Theory (M1, M2)\\
    $\bullet$ Optimal Transport (M3)\\
    $\bullet$ State Space Models (M5)\\
    $\bullet$ Bayesian Inference (M6)\\
    $\bullet$ Game Theory (M7)
};

\end{tikzpicture}
\caption{Hierarchical decomposition of two complementary master problems into 14 sub-problems addressed by 7 novel methods. Master Problem A focuses on adversarial resilience across heterogeneous networks, while Master Problem B optimizes fundamental trade-offs between robustness, accuracy, and privacy. Each method addresses multiple sub-problems through integrated mathematical frameworks including neural ODEs, optimal transport, graph theory, state space models, Bayesian inference, and game theory. Performance metrics shown represent key achievements on benchmark datasets.}
\label{fig:problem_hierarchy}
\end{figure}
